\clearpage
\section{Anexo 8.I: Observaciones, equivalencias e integración externa}
\label{sec:cierre-traza}
El correo principal, punto 1.b, pide una ``Tabla de trazabilidad observación--\allowbreak respuesta--\allowbreak sección modificada'' de las observaciones del Informe 2. Se aplica al alcance S8. Las citas textuales siguientes proceden de la revisión G02--LafroX, pp. 33--38; Bases Administrativas, apartado Subdocumento 8, p. impresa 58, confirma técnica, costo-beneficio y reservas. Correo principal prevalece sobre las Bases cuando las aclara; no contiene una nueva lista técnica particular de riesgos. Los estados distinguen cumplimiento documental en S8 de integración o ejecución externa. No se atribuye a una tabla prevista valor de prueba realizada.

\begingroup
\renewcommand{\thetable}{8.I.1}
\begin{landscape}
\setlength{\textwidth}{\linewidth}
\renewcommand{\arraystretch}{1.08}
\begin{tablalafrox}{Trazabilidad observación--respuesta--sección modificada. Fuente: correo principal 1.b y revisión Informe 2, pp. 33--38.}{tab:cierre-observaciones}{G{0.05} F{0.28} Y F{0.15} F{0.15}}{\cab{Obs.} & \cab{Pedido original y fuente} & \cab{Respuesta aplicada} & \cab{Ubicación} & \cab{Estado}}
01 & ``Definir los factores decimales de impacto en el cuerpo, corregir el solape de bandas, renumerar tablas'' (p. 38). & Factores, bandas y umbrales únicos; escala de HH condicional y máximo de dimensiones. & 8.1.3; tablas 8.2--8.4; 8.C.1. & corregido en S8 \\
02 & ``Para cada riesgo que depende del CLIENTE, nombrar la contraparte, qué se le pide y en qué fecha, y la contingencia si no responde'' (p. 37). & Se conservan 28 coordinaciones, plazos y escalamiento; la ficha agrega esfuerzo y plazo de contingencia. & 8.A.2; 8.F.1; T-16. & corregido en S8 \\
03 & ``Retirar del registro los 13 riesgos internos del proveedor [...] y resolverlos en la propuesta'' (p. 37). & Treinta y cinco riesgos vigentes; 16 CI con solución y cierre propuestos, incluidos los trece observados. Incorporar en S4/S9/S13 y verificar; no se declara ejecutado. & 8.E.1; equivalencias 8.I.2. & requiere integración externa \\
04 & ``Declarar el riesgo residual de cada riesgo después de la respuesta'' (p. 37). & 35 valoraciones posteriores proyectadas con probabilidad, impacto, zona, hipótesis y evidencia. Eficacia efectiva se verifica al ejecutar. & 8.F.2; resumen posterior en T-16. & corregido en S8 \\
05 & ``ningún riesgo puede quedar en Alta o Crítica sin un plan adicional'' (p. 37). & Los que mantienen E alta/crítica reciben plan adicional, plazo, dueño y bloqueo; no se acredita todavía un nivel manejable ni se rebaja I sanitario o contractual sin fundamento. & 8.F.1--2; 8.G.1. & parcial \\
06 & ``Agregar los riesgos del Cap. 19 que faltan [...] y los de obsolescencia, bloqueo por proveedor y escalabilidad'' (p. 37). & Se conservan omisiones incorporadas, sus fuentes y no duplicación. Escalabilidad y consistencia del diseño se resuelven como CI; cambios externos de soporte/salida siguen como riesgos. & 8.D.3; 8.E.1; T-16; 8.I.2. & parcial \\
07 & ``Llevar al registro los puntos únicos de falla y riesgos residuales del S4'' (p. 37). & ERP externo sigue en T-16; identidad y stock quedan como controles internos trazados, no se trasladan defectos al CLIENTE. Faltan evidencia y sincronización de S4/S9. & 8.D.1--2; 8.E.1; 8.F.2. & requiere integración externa \\
08 & ``Aplicar una técnica de las exigidas (FMEA, árbol de fallas o simulación) a los riesgos críticos y altos'' (p. 38). & FMEA de 15 modos, causas/efectos, S/O/D, NPR, controles, acciones y puntuación posterior proyectada. & 8.G.1; síntesis 8.2.3. & corregido en S8 \\
09 & ``Hacer el costo-beneficio de cada respuesta de zona alta o crítica, con plazo, y justificar las de retorno menor que 1'' (p. 38). & Quince cálculos en HH, entradas A/C/p0/p1/H0, sensibilidad y justificación obligatoria; calendario en ficha. & 8.F.1/8.F.3. & corregido en S8 \\
10 & ``ninguna acción tiene plazo ni fecha [...] ni esfuerzo, ni contingencia específica'' (p. 36). & Treinta y cinco fichas con perfiles, HH extra y actuación contingente temporizada; controles ordinarios permanecen en EDT. & 8.F.1; 8.A.1--2. & corregido en S8 \\
11 & ``Derivar la reserva de gestión, conciliar contingencia y reserva protegida del S7, e incorporarlas al T-15 y al cronograma'' (p. 38). & Gestión propuesta derivada por escenario, conciliación sin doble conteo y distribución mensual/perfil; integración de cronograma/flujo pendiente. & 8.H.1--3; 8.3.2. & requiere integración externa \\
12 & ``El S13 dice que el Anexo 8.A asigna al piloto de R-034 [histórico] [...] meses 11 a 15'' (p. 36). & CI-034: paquete 11--15, certificación 11--12, piloto real 13--15; conservar indicador y actualizar referencia S13. & 8.D.1/TR-01; 8.E.1. & requiere integración externa \\
13 & ``Dos tablas distintas numeradas Tabla 8.4'' y ``La Tabla 8.5 tiene siete columnas''; también se observa la fuente de RT-20.04/RT-20.02 y la ausencia de folio/firma en RBS (p. 37). & Numeración única, tablas de resumen de cinco columnas y referencias BTT cap. 20 corregidas. RBS conserva folios y firma. & 8.1.3; 8.2; 8.3.3. & corregido en S8 \\
14 & ``quitar académico'' (p. 38). & No se aplica por instrucción expresa del autor de conservar las citas que contienen contenido academico; se registra la excepción, no se declara cumplido. & Tablas 8.2--8.4 y declaración de IA protegidas. & parcial \\
\end{tablalafrox}
\end{landscape}
\endgroup

\begingroup
\renewcommand{\thetable}{8.I.2}
\begin{landscape}
\setlength{\textwidth}{\linewidth}
\renewcommand{\arraystretch}{1.08}
\begin{tablalafrox}{Equivalencias completas: ID anterior histórico y destino vigente. Fuente: catálogo anterior de 51 eventos; nuevo T-16 de 35 riesgos.}{tab:cierre-equivalencias}{F{0.14} G{0.11} Y F{0.20}}{\cab{ID anterior} & \cab{ID vigente / CI} & \cab{Descripción} & \cab{Estado}}
Histórico R-001 & R-001 & Cadena no completa a tiempo su certificación EDI & riesgo vigente \\
Histórico R-002 & R-002 & Historia sanitaria de origen no reconstruible & riesgo vigente \\
Histórico R-003 & R-003 & Historial de ventas inutilizable para carga y reposición & riesgo vigente \\
Histórico R-004 & R-004 & Diferencias de inventario sin resolver antes del corte & riesgo vigente \\
Histórico R-005 & R-005 & Saldo inicial de envases no conciliable (carga M8) & riesgo vigente \\
Histórico R-006 & R-006 & Historia de origen insuficiente para líneas base de aceptación & riesgo vigente \\
Histórico R-007 & R-007 & Deuda viva adicional a la cartera de la volumetría & riesgo vigente \\
Histórico R-008 & R-008 & SII no responde y deja un trámite tributario pendiente & riesgo vigente \\
Histórico R-009 & CI-009 & Dado que la integración de pagos depende de la respuesta del POS de un tercero y un timeout no permite concluir que el cargo fue rechazado, entonces puede que el POS agote el tiempo de respuesta sin permitir comprobar si se efectuó el cargo, teniendo como consecuencia que el estado del pago quede incierto y sea necesario investigar y conciliarlo con la contraparte financiera antes de resolver la incidencia. & control interno; 8.E.1 \\
Histórico R-010 & CI-010 & Dado que la consulta de ubicaciones y la planificación de rutas dependen de un servicio externo de mapas, entonces puede que el servicio externo deje de estar disponible durante la operación, teniendo como consecuencia que se retrasen o interrumpan las tareas de planificación que dependen del servicio y deban atenderse la incidencia y la recuperación. & control interno; 8.E.1 \\
Histórico R-011 & R-009 & Calidad no entrega la matriz térmica aprobada & riesgo vigente \\
Histórico R-012 & R-010 & Equipos del cliente no llegan antes de la ola & riesgo vigente \\
Histórico R-013 & R-011 & Transportista no firma a tiempo el acuerdo de incorporación & riesgo vigente \\
Histórico R-014 & R-012 & Transportista no permite instalar el termógrafo & riesgo vigente \\
Histórico R-015 & R-013 & Usuarios mantienen el proceso anterior & riesgo vigente \\
Histórico R-016 & R-014 & Planificador no disponible antes de validar excepciones & riesgo vigente \\
Histórico R-017 & R-015 & Usuarios no asisten a la capacitación antes de la ola & riesgo vigente \\
Histórico R-018 & R-016 & Contraparte no entrega validación en la fecha prevista & riesgo vigente \\
Histórico R-019 & R-017 & Proveedor de conectividad interrumpe el enlace de un sitio & riesgo vigente \\
Histórico R-020 & R-018 & Corte eléctrico supera el tiempo de respaldo & riesgo vigente \\
Histórico R-021 & R-019 & Falta aprobación del cliente para datos en regiones de nube & riesgo vigente \\
Histórico R-022 & R-020 & Cliente no autoriza intervención urgente en congelamiento & riesgo vigente \\
Histórico R-023 & CI-023 & Dado que las credenciales pueden ser obtenidas mediante engaño, phishing u otros métodos y usadas sin autorización, entonces puede que un tercero utilice credenciales ajenas para autenticarse como otra persona e intente registrar pedidos o acceder a portales, teniendo como consecuencia que deban investigarse y contenerse los accesos u operaciones afectados y, si la suplantación tiene éxito, puedan atribuirse acciones al titular legítimo. & control interno; 8.E.1 \\
Histórico R-024 & CI-024 & Dado que la solución recibe de cadenas comerciales mensajes electrónicos de pedidos y avisos de despacho mediante sistemas externos, entonces puede que llegue un mensaje que aparente provenir de una cadena autorizada sin serlo, o cuyo contenido haya sido alterado, teniendo como consecuencia que deba aislarse e investigarse el mensaje antes de procesar operaciones y, si se acepta, pueda alterarse el pedido o aviso correspondiente. & control interno; 8.E.1 \\
Histórico R-025 & CI-025 & Dado que los terminales enrolados se utilizan físicamente en ruta y bodega y contienen datos locales asociados a usuarios autorizados, entonces puede que un terminal enrolado se pierda o quede en poder de una persona distinta de su usuario autorizado, teniendo como consecuencia que deban investigarse la posible exposición de datos locales y contenerse accesos indebidos al dispositivo. & control interno; 8.E.1 \\
Histórico R-026 & CI-026 & Dado que los portales e interfaces públicas están expuestos a solicitudes de actores externos, entonces puede que un actor dirija ráfagas abusivas de solicitudes para impedir accesos legítimos, teniendo como consecuencia que se degraden los accesos públicos y se requiera contención y verificación de la disponibilidad del servicio. & control interno; 8.E.1 \\
Histórico R-027 & CI-027 & Dado que las cuentas administrativas están expuestas a campañas de phishing y pueden dar acceso privilegiado a la plataforma, entonces puede que un atacante obtenga credenciales administrativas y las utilice para borrar o alterar respaldos o evidencia, teniendo como consecuencia que deban contenerse los accesos comprometidos e investigarse y verificarse la integridad y recuperabilidad de los conjuntos afectados. & control interno; 8.E.1 \\
Histórico R-028 & R-021 & Fallas mecánicas del camión & riesgo vigente \\
Histórico R-029 & R-022 & Camión involucrado en un siniestro & riesgo vigente \\
Histórico R-030 & R-023 & Asalto a camión & riesgo vigente \\
Histórico R-031 & R-024 & Reemplazos de conductores externos sin aviso & riesgo vigente \\
Histórico R-032 & R-025 & Solicitudes que amplían el alcance & riesgo vigente \\
Histórico R-033 & R-026 & Cambios normativos & riesgo vigente \\
Histórico R-034 & CI-034 & Dado que el aviso de vencimiento posentrega de la innovación 1 estima el ritmo de venta del local a partir de las compras del canal tradicional, que son irregulares, entonces puede que los avisos resulten falsos o no tengan una acción posible por falta de catálogo, teniendo como consecuencia que el preventista y el almacenero dejen de confiar en ellos y que la función no alcance el indicador I-01A antes de liberarse. & control interno; 8.E.1 \\
Histórico R-035 & CI-035 & Dado que la reproducción de incidentes de terreno de la innovación 2 reconstruye el estado del terminal a partir de paquetes capturados en operación, entonces puede que la reproducción no explique el incidente observado o que la captura consuma batería o datos del terminal, teniendo como consecuencia que soporte deba volver al diagnóstico convencional y que no se obtenga la reducción del tiempo de diagnóstico del indicador I-02A. & control interno; 8.E.1 \\
Histórico R-036 & CI-036 & Dado que los paquetes de reproducción de la innovación 2 se capturan en terminales que operan con datos de clientes y conductores, entonces puede que un paquete conserve datos personales sin sanitizar, teniendo como consecuencia una exposición de datos personales y la suspensión del conjunto de paquetes afectado. & control interno; 8.E.1 \\
Histórico R-037 & CI-037 & Dado que la vida útil remanente de la innovación 3 se estima con parámetros por familia de producto a partir de la historia térmica del lote, entonces puede que los parámetros no representen los productos de Puelche o que Operaciones lea la estimación como permiso para relajar la regla RNG-04, teniendo como consecuencia órdenes de salida sugeridas erróneas para productos refrigerados y congelados. & control interno; 8.E.1 \\
Histórico R-038 & CI-038 & Dado que la estimación de la innovación 3 asocia cada lote con las lecturas de los sensores por los que pasó, entonces puede que un lote se mueva sin registrar su evento de ubicación y quede asociado a un sensor equivocado, teniendo como consecuencia una historia térmica errónea del lote y una estimación de vida útil sin fundamento. & control interno; 8.E.1 \\
Histórico R-039 & R-027 & Disputa en la liquidación del tramo variable & riesgo vigente \\
Histórico R-040 & CI-040 & Dado que la hoja de negocio de la innovación 5 entrega al almacenero indicadores de su propio local y, desde el bloque 4, comparaciones agregadas, entonces puede que el almacenero no la lea o no la entienda, o que de un bloque comparativo se puedan inferir cifras de un vecino, teniendo como consecuencia que no se alcance el indicador I-05B o que deba desactivarse el bloque 4. & control interno; 8.E.1 \\
Histórico R-041 & R-028 & Desacuerdo sobre telemetría impide incorporación & riesgo vigente \\
Histórico R-042 & R-029 & Rotación nocturna supera inducción disponible & riesgo vigente \\
Histórico R-043 & R-030 & Interfaces ERP/WMS no verificables a tiempo & riesgo vigente \\
Histórico R-044 & R-031 & Peak y congelamiento impiden cierre de fase & riesgo vigente \\
Histórico R-045 & R-032 & Cambio de soporte exige adecuación extraordinaria & riesgo vigente \\
Histórico R-046 & R-033 & Dependencia de proveedor limita salida o sustitución & riesgo vigente \\
Histórico R-047 & CI-047 & Dado que el peak alcanza unas 2.600 entregas diarias y concentra despacho y reconexión, entonces puede que una dependencia o perfil de proceso no sostenga la carga efectiva aun con el escalado previsto, teniendo como consecuencia colas crecientes, incumplimiento de respuesta o indisponibilidad de operaciones críticas. & control interno; 8.E.1 \\
Histórico R-048 & CI-048 & Dado que altas, bajas y relevos dependen de identidad central y permisos locales vigentes, entonces puede que la indisponibilidad del IdP o la falta de un manifiesto válido impida habilitar al personal requerido o propagar una revocación, teniendo como consecuencia afectación de continuidad y seguridad del relevo. & control interno; 8.E.1 \\
Histórico R-049 & R-034 & ERP indisponible impide nueva guía válida & riesgo vigente \\
Histórico R-050 & CI-050 & Dado que la reserva global depende del estado central y de una única autoridad de escritura, entonces puede que falle la consolidación o no pueda restituirse una autoridad consistente, teniendo como consecuencia stock central desactualizado y pedidos sin confirmación o con conflictos de reserva. & control interno; 8.E.1 \\
Histórico R-051 & R-035 & Conductor externo rechaza uso de la solución & riesgo vigente \\
\end{tablalafrox}
\end{landscape}
\endgroup
\paragraph{Regla de identificación.} R-001 a R-035 son exclusivamente el catálogo vigente. Las columnas rotuladas ``Histórico'' contienen los ID anteriores y no deben usarse como claves de T-16. Un CI mantiene su sufijo histórico, sin ocupar un número R. Al sincronizar otros documentos se consulta esta tabla; no se realiza sustitución global de dígitos, porque un ID antiguo puede coincidir con otro vigente. Las citas originales de la revisión conservan su ID histórico rotulado, sin alterar el contenido citado. Los expedientes futuros usarán ID vigente y referencia histórica sólo en metadatos.

\paragraph{Cambios externos concretos.} En S4, enlazar CI-009/010 y CI-023--027 con integración, mapas y seguridad, y CI-047/048/050 con capacidad, identidad y reconciliación; incorporar soluciones y casos de 8.E, conservar las dependencias externas del nuevo T-16 mediante equivalencia. En S9/T-13, registrar casos, requisitos, versión y evidencia de esos controles, FMEA y bloqueo habilitante; la tabla prevista no equivale a ensayo ejecutado. En S13, convertir antiguos riesgos de diseño INN-01/02/03/05 en CI y criterios de validación; mantener CI-034 y su piloto 13--15 dentro del paquete 11--15, y enlazar liquidación variable con el riesgo vigente equivalente al histórico R-039. En S7/T-14/T-15, cargar la propuesta de 8.H sólo después de conciliar paquetes y capacidad y conservar protección E1; entregar el mismo calendario al modelo económico sin duplicar reserva. En S6 y matrices transversales, actualizar identificadores con 8.I.2 y responsables/decisiones. JP registra aceptación de cada sincronización; hasta editar y comprobar esos documentos el estado sigue siendo requiere integración externa.
